\def\paperid{A08}
\documentclass[10pt,a4paper]{article}
\usepackage[margin=23mm,top=24mm,bottom=24mm]{geometry}
\usepackage{fontspec,polyglossia,amsmath,amssymb,booktabs,array,enumitem,xcolor,fancyhdr,xurl,hyperref,microtype}
\setdefaultlanguage{spanish}\setmainfont{Latin Modern Roman}\setsansfont{DejaVu Sans}
\definecolor{ink}{HTML}{17375C}\definecolor{mint}{HTML}{0F8F86}
\hypersetup{colorlinks=true,urlcolor=mint,linkcolor=ink,pdfauthor={Juan Manuel Goenaga Castro}}
\setlength{\parindent}{0pt}\setlength{\parskip}{5pt}\linespread{1.05}\emergencystretch=2em
\setlength{\headheight}{15pt}\pagestyle{fancy}\fancyhf{}
\fancyhead[L]{\sffamily\footnotesize GoToGym / Quantum Coach}\fancyhead[R]{\sffamily\footnotesize Robustecimiento QAF / \paperid}
\fancyfoot[L]{\scriptsize Prepublicación metodológica / 22 de octubre de 2026}\fancyfoot[R]{\small\thepage}
\setlist{nosep,leftmargin=1.5em}
\begin{document}
\begin{center}{\sffamily\small\textcolor{mint}{SERIE DE ROBUSTECIMIENTO / A08}}\par\vspace{10pt}{\sffamily\LARGE\bfseries\textcolor{ink}{Operación reproducible y entrega verificable}}\par\vspace{7pt}{\large Caché por contenido, salida atómica y estudio integrado}\par\vspace{12pt}\textbf{Juan Manuel Goenaga Castro}\par GoToGym Quantum Research\par 22 de octubre de 2026 / versión 1.1\end{center}\colorbox{ink!6}{\parbox{.96\linewidth}{\small\textbf{Estado.} Prepublicación metodológica sin revisión por pares. Propuesta computacional implementada; evaluación perceptual y estudio con participantes pendientes.}}\begin{abstract}Se integran compositor, storyboard, voz, pantalla y auditoría. Se implementa una clave con contenido y entorno, salida temporal, reemplazo atómico tras verificación, una generación de video y decodificación completa. QAF aporta coordinación de estados y restricciones; no se atribuye una mejora de productividad sin observar varios productores. La instalación de modelos sigue siendo un paso opcional separado. Se presenta el problema observado, una formalización con unidades y restricciones, el código correspondiente y sus controles. Se distingue la verificación de ingeniería de calidad percibida, identidad regional y efecto comercial. El uso QAF corresponde a la región clásica identificable y conserva un comparador con los mismos recursos.\end{abstract}
\section{Diagnóstico y pregunta de investigación}
El episodio funcionaba mediante programas externos al estudio, dependencias implícitas y una segunda codificación para quemar subtítulos. Un archivo corrupto podía parecer reutilizable por su mera presencia. Un mensaje de trabajo terminado podía interpretarse como aprobación humana. Se necesitaba un proceso instalable y verificable desde la misma plataforma.

Se integran compositor, storyboard, voz, pantalla y auditoría. Se implementa una clave con contenido y entorno, salida temporal, reemplazo atómico tras verificación, una generación de video y decodificación completa. QAF aporta coordinación de estados y restricciones; no se atribuye una mejora de productividad sin observar varios productores. La instalación de modelos sigue siendo un paso opcional separado.
\section{Modelo matemático y restricciones}
Se registran proyecto canónico, hashes de medios y fuentes, versión del motor, subtítulos, formato y versiones de Python, Pillow y FFmpeg. La salida se comprueba por cuadros decodificados, duración, tamaño y audio. Una clave se compara con otra solo cuando la codificación y dependencias son compatibles.
\begin{equation}K=H(\operatorname{canon}(P,A,C,V,E,F,CC)).\end{equation}

\begin{equation}G_{\rm cache}=\mathbf1[K=K_{\rm anterior}]\mathbf1[H(M)=H_{\rm registrado}].\end{equation}

\begin{equation}G_{\rm entrega}=G_{\rm decode}\land G_{\rm cuadros}\land G_{\rm audio}\land G_{\rm tamano}.\end{equation}
La serialización canónica elimina diferencias irrelevantes de orden JSON. Una clave coincidente no basta si el MP4 está truncado; la identidad de bytes detecta ese caso. La publicación por rename ocurre después de decodificar y medir, dentro del mismo sistema de archivos. Los subtítulos se dibujan en los cuadros de entrada, evitando una generación adicional con pérdida. El coste total distingue trabajo y tiempo de pared. No se compara velocidad entre máquinas sin describir hardware y carga.
\subsection{Región identificable QAF y comparador}
La integración usa el caso diagonal del formalismo, con base fija de candidatos y compuerta de factibilidad g:
\begin{equation}
p_j=\frac{g_j\exp[-(L_j-L_{\min})/\tau]}{\sum_k g_k\exp[-(L_k-L_{\min})/\tau]},\quad\rho=\operatorname{diag}(p),\quad\rho\succeq0,\quad\operatorname{Tr}\rho=1.
\end{equation}
Estos pesos describen decisión, no satisfacción del público. No se aprende ni asigna fase sin observaciones que la identifiquen. La enumeración nativa selecciona el mínimo factible; el comparador clásico recibe exactamente los mismos candidatos y coincide por construcción. No se reclama ventaja cuántica. MPS y TDVP siguen disponibles como referencias matemáticas anteriores, pero no intervienen en la síntesis de voz, el transporte de luz o una supuesta dinámica cerebral.

Con pérdidas desplazadas no negativas, se conserva la compilación triangular de selección única:
\begin{equation}
E(x)=\sum_j(\widetilde L_j-A)x_j+2A\sum_{j<k}x_jx_k+A,\qquad A>\max_j\widetilde L_j.
\end{equation}
La constante distingue correctamente el conjunto vacío; la interacción aparece una sola vez por pareja. La máscara de alternativas prohibidas es una restricción nativa, no una penalización blanda. Una prueba independiente enumera las 64 configuraciones de seis variables y compara con el óptimo restringido. La positividad y la traza se verifican por separado. Este uso explícito de la región clásica evita forzar complejidad tensorial sin evidencia incremental.
\section{Implementación y trazabilidad}
cinema\_\allowbreak{}pipeline.py implementa render, auditoría, storyboard y caché. Los comandos cinema-audit, cinema-storyboard y cinema-render comparten motor con el servidor. La interfaz ofrece pantalla original y tasa objetivo; los lanzadores abren el estudio local. ASR puede operar sin cargar pesos TTS. Los perfiles finales exigen firmas humanas documentadas y nunca las fabrican a partir de un resultado técnico.

Módulos: \texttt{cinema\_\allowbreak{}pipeline.py}, \texttt{server.py}, \texttt{cli.py}, \texttt{release.py}. Cada transformación conserva archivos de entrada, parámetros y hashes. No se sustituye una capacidad ausente por una declaración de disponibilidad.
\subsection{Controles que pueden refutar la implementación}
\begin{flushleft}\small test\_\allowbreak{}both\_\allowbreak{}aspects\_\allowbreak{}single\_\allowbreak{}encode\_\allowbreak{}full\_\allowbreak{}decode\_\allowbreak{}and\_\allowbreak{}verified\_\allowbreak{}cache\\ test\_\allowbreak{}cache\_\allowbreak{}invalidates\_\allowbreak{}captions\_\allowbreak{}fonts\_\allowbreak{}engine\_\allowbreak{}environment\\ suite del servidor\\ dos exportaciones completas.\end{flushleft}

La batería combina invariantes algebraicos, casos adversos y render pequeño con decodificación. El conjunto de tests es un piso de reproducibilidad: aprobar una identidad de bytes no acredita veracidad de contenido, y una ecuación correcta no acredita preferencia artística.
\section{Resultados computacionales observados}
Se ejecutaron pruebas pequeñas en ambos formatos, reutilización con hash válido y reconstrucción de un MP4 deliberadamente corrupto. Las dos exportaciones completas pasaron decodificación, cuadros, duración, tamaño y audio. El compositor genera video una sola vez por archivo y comparte implementación con CLI y servidor.

La suite ejecutada contiene 72 pruebas aprobadas. Son contratos de software e integración, no 72 observaciones independientes de calidad artística. Los datos completos se registran en el informe de versión.

El archivo \texttt{docs/validation/robustness-v0.4.json} conserva medidas y hashes de entrega. El proyecto completo con voz se distribuye por separado del código público. No se publican notas originales, perfiles biométricos ni sus transcripciones privadas.
\section{Protocolo de evaluación y ablación}
Aplicar cambios de texto, fuentes, captura, códec y versiones; medir invalidaciones correctas y recuperación tras corrupción. Comparar proceso externo y estudio integrado con iguales recursos, operadores e intentos. Medir tiempo hasta entrega repetible y defectos escapados. Pedir instalación en un equipo distinto sin rutas absolutas del autor. Registrar fallos y dependencias opcionales; la coincidencia de entradas no garantiza determinismo universal de GPU o modelos neuronales.

La aceptación es jerárquica: primero se satisfacen las restricciones duras; después se optimiza calidad dentro del conjunto admisible. Se registrarán candidatos fallidos, tiempo y recursos, no solo la toma ganadora. El episodio, la persona y la familiaridad se controlan por separado. Repetir valoraciones de un único video no crea videos independientes. La muestra confirmatoria se calculará con variabilidad del piloto, sin inventar potencia o resultados.
\section{Discusión, límites y propiedad intelectual}
La clave incluye todo el código del paquete y puede invalidar más de lo estrictamente necesario. No se ha construido una granja distribuida ni un servicio multiusuario de internet. La película es un animatic fotográfico. Preferencia artística, naturalidad, acento y efecto comercial conservan protocolos independientes.

Los aportes de integración se preparan a nombre de Juan Manuel Goenaga Castro. Se conserva atribución de Antoni Bancells al antecedente técnico QAF y de los modelos, métodos y fuentes de terceros. No se reivindica entrenamiento propio de Qwen o Whisper. Una ampliación relacional deberá superar un comparador clásico con las mismas entradas y presupuesto; si no aporta valor, la solución sencilla continúa siendo válida.
\section{Conclusión}
La propuesta computacional descrita se encuentra en el código entregado y permite repetir sus controles. La calidad percibida, la autenticidad vocal y la eficacia comercial requieren sus propios datos; no se declaran resueltas por completar una exportación.
\begin{thebibliography}{9}
\bibitem{qaf} Goenaga Castro, J. M. (2025--2026). Quantum Analysis Frameworks y serie QAF audiovisual I--V, documentos aportados. Bancells, A. (2026). Dinámica operacional, MPS/TDVP y compilación QUBO, antecedente técnico aportado. Atribuciones conservadas; sus bancos completos no se repitieron en esta ampliación.
\bibitem{ffmpeg} FFmpeg. Documentación primaria de filtros atempo, loudnorm, sidechaincompress y scale. Consulta 22/10/2026. \url{https://ffmpeg.org/ffmpeg-filters.html}.
\end{thebibliography}
\end{document}
